\subsection{诺特环上系数的多项式}

设 A 为环，\(\sigma\) 为环 A 的自同态，d 为 A 的加群的自同态，满足关系

\[
d (a b) = \sigma (a) d (b) + d (a) b\tag{1}
\]

其中 \(a, b \in A\) 任意。换句话说，d 是从环 A 到 (A,A)-双模的一个导子，该双模是在 A 的加群上赋予左作用律 \((a, x) \mapsto \sigma(a)x\) 和右作用律 \((x, a) \mapsto xa\) 而得到的。我们有 \(d(1) = 0\) （III, §10, No.~5, p.~557, 命题 3）。

回顾（IV, §1, No.~1, p.~2）A{[}X{]} 表示系数在 A 中的一元多项式组成的 Z-模 \(A \otimes_{Z} Z[X]\) 。我们赋予它左 A-模的自然结构。族 \((\mathrm{X}^{n})_{n \in \mathbf{N}}\) 是 A{[}X{]} 在 A 上的一个基。我们把 A 与它在映射 \(a \mapsto a \otimes 1\) 下的像等同起来。

命题 7. --- 设 A、\(\sigma\) 、\(d\) 如上。在群 A{[}X{]} 上存在唯一的环结构，具有下列性质：

\begin{enumerate}
\def\labelenumi{\alph{enumi})}
\item
  此环中的加法是 A{[}X{]} 的通常加法。
\item
  此环中的乘法延拓 A 的乘法。
\item
  在此环中，由 A 的一个元素 a 后接 n 个等于 X 的项组成的序列 \((a,X,\ldots,X)\) 的乘积是多项式 \(aX^{n}\) 。
\item
  在此环中，\(\mathrm{X}a = \sigma(a)\mathrm{X} + d(a)\) 对每个 \(a \in \mathrm{A}\) 成立。
\end{enumerate}

设 E 为加群 A{[}X{]} 的自同态环。把 \(a \in A\) 送到位似 \(a_{M}\) ——即左 A-模 M = A{[}X{]} 的位似——的映射，是从 A 到 E 的环同态。考虑 E 的元素 u、\(\sigma_{M}\) 与 \(d_{M}\) ，它们由 \(u(\sum b_{n}X^{n}) = \sum b_{n}X^{n+1}\) ，\(\sigma_{\mathrm{M}}(\sum b_{n}X^{n}) = \sum \sigma(b_{n})X^{n}\) ，\(d_{\mathrm{M}}(\sum b_{n}X^{n}) = \sum d(b_{n})X^{n}\) 定义。对每个 \(a \in A\) ，我们有

\[
u a _ {\mathrm{M}} = a _ {\mathrm{M}} u, \quad \sigma_ {\mathrm{M}} a _ {\mathrm{M}} = \sigma (a) _ {\mathrm{M}} \sigma_ {\mathrm{M}}, \quad d _ {\mathrm{M}} a _ {\mathrm{M}} = \sigma (a) _ {\mathrm{M}} d _ {\mathrm{M}} + (d (a) _ {\mathrm{M}}).\tag{2}
\]

\[
X_ {M} = \sigma_ {M} u+ d_ {M} .\tag{3}
\]

由 (2) 推出，对每个 \(a \in \mathbf{A}\) ，有

\[
\mathrm{X} _ {\mathrm{M}} a _ {\mathrm{M}} = \sigma (a) _ {\mathrm{M}} \mathrm{X} _ {\mathrm{M}} + (d (a) _ {\mathrm{M}}).\tag{4}
\]

考虑映射 \(\varphi : \mathrm{A}[\mathrm{X}] \to \mathrm{E}\) ，它由

\[
\varphi \Bigl (\sum a _ {n} \mathrm{X} ^ {n} \Bigr) = \sum (a _ {n}) _ {\mathrm{M}} (\mathrm{X} _ {\mathrm{M}}) ^ {n}.\tag{5}
\]

定义。这是一个群同态。一个归纳论证表明，\((\mathrm{X}_{\mathrm{M}})^{n}(1)=\mathrm{X}^{n}\) 对每个 \(n\in N\) 成立。于是 \(\varphi(\mathrm{P})(1)=\mathrm{P}\) 对每个 \(P\in\mathrm{A}[X]\) 成立，这就证明同态 \(\varphi\) 是单的。我们用 B 表示它的像。集合 B 是 E 的子群；它包含 1，并且对 \(a_{M}\) （\(a\in A\)）和 \(X_{M}\) 的左乘法封闭（见 (4)）。因此它是 E 的子环。A{[}X{]} 上由经 \(\varphi\) 的结构转移从 B 的环结构导出的唯一环结构具有命题 7 的性质，其中性质 d) 由关系 (4) 得出。

若 A{[}X{]} 带有一个具有命题 7 的性质的环结构，则此环的左位似 \(\gamma_{X}\) （I, §8, No.~1, p.~97）必然把 \(bX^{n}\) 送到 \(\sigma(b)\mathrm{X}^{n+1} + d(b)\mathrm{X}^{n}\) （对 \(b \in A\) 和 \(n \in N\) ），故等于 \(X_{M}\) 。位似 \(\gamma_{a}\) 必然等于 \(a_{M}\) ，对每个 \(a \in A\) 。于是，\(\gamma_{\mathrm{P}} = \varphi(\mathrm{P})\) 对每个 \(P \in A[X]\) 成立；命题 7 中的唯一性由此得出。

带有具有命题 7 的性质的唯一环结构的集合 A{[}X{]} 记作 \(A[X]_{\sigma,d}\) ，称为 X 的、系数在 A 中且相对于 \(\sigma\) 与 d 的多项式环。~当 d 是零映射时，我们简记为 \(A[X]_{\sigma}\) ；记号 A{[}X{]} 则用于进而 \(\sigma\) 是 A 上的恒等映射的情形。这一记号与 IV, §1, No.~1, p.~1 对交换环 A 引入的记号相容。

注. --- 环 \(A[X]_{\sigma,d}\) 具有下列泛性质：给定环 \(A'\) 、环同态 \(f: A \to A'\) 以及 \(A'\) 的元素 x，使得 \(xf(a) = f(\sigma(a))x + f(d(a))\) 对每个 \(a \in A\) 成立，则存在唯一的环同态 \(g: A[X]_{\sigma,d} \to A'\) ，它延拓 f 且把 X 映到 x。

唯一性是明显的。因此我们来证明映射 \(g: \mathrm{A}[\mathrm{X}]_{\sigma,d} \to \mathrm{A}'\) （它由 \(g(\sum a_n\mathrm{X}^n) = \sum f(a_n)x^n\) 定义）具有所需的性质。它延拓 \(f\) ，把 \(\mathrm{X}\) 映到 \(x\) ，且是群同态。我们有 \(g(1) = 1\) 。对 \(a \in \mathrm{A}\) 以及 \(\mathrm{Q} = \sum a_n\mathrm{X}^n\) （其中 Q 属于 \(\mathrm{A}[\mathrm{X}]_{\sigma,d}\) ），有

\[
g (a \mathrm{Q}) = g \left(\sum a a _ {n} X ^ {n}\right) = \sum f (a a _ {n}) x ^ {n} = f (a) \sum f (a _ {n}) x ^ {n} = g (a) g (\mathrm{Q})
\]

以及

\[
\begin{array}{l} g (\mathrm{XQ}) = g \Bigl (\sum \Bigl (\sigma (a _ {n}) \mathrm{X} ^ {n + 1} + d (a _ {n}) \mathrm{X} ^ {n} \Bigr) \Bigr) \\ = \sum (f (\sigma (a _ {n})) x ^ {n + 1} + f (d (a _ {n})) x ^ {n}) = x \sum f (a _ {n}) x ^ {n} = g (\mathrm{X}) g (\mathrm{Q}). \end{array}
\]

由此推出，\(g(\mathrm{P})g(\mathrm{Q}) = g(\mathrm{PQ})\) 对 \(A[X]_{\sigma,d}\) 中的 P、Q 成立，因而 g 是环同态。

定理 2. --- 设 A 是左诺特环，\(\sigma\) 是 A 的一个自同构，\(d\) 是 A 的加法群的一个满足关系式 (1) 的自同态。则环 \(\mathrm{A}[\mathrm{X}]_{\sigma,d}\) 是左诺特的。

令 \(B = A[X]_{\sigma,d}\) 。对任意整数 \(n \geqslant 0\) ，我们用 \(B_{n}\) 表示 B 中形如 \(a_{0} + a_{1}X + \cdots + a_{n}X^{n}\) 的元素组成的集合。它是 B 的一个左 A-子模。映射 \(\varphi_{n}: B_{n} \to A_{s}\) （由 \(\varphi_{n}(a_{0} + a_{1}X + \cdots + a_{n}X^{n}) = a_{n}\) 定义）是 A-线性的。

设 \(\mathfrak{b}\) 是 B 的一个左理想。对每个整数 \(n \geqslant 0\) ，集合 \(\mathfrak{a}_n = \varphi_n(\mathfrak{b} \cap \mathrm{B}_n)\) 是 A 的一个左理想。由于 \(Xa = \sigma(a)X + d(a)\) 对每个 \(a \in A\) 成立，我们有

\[
\varphi_ {n + 1} (\mathrm{XQ}) = \sigma (\varphi_ {n} (\mathrm{Q}))\tag{6}
\]

对每个 \(Q \in B_{n}\) 成立，因而 \(\sigma(\mathfrak{a}_{n}) \subset \mathfrak{a}_{n+1}\) 。于是，A 的理想序列 \(\mathfrak{a}_{n}^{\prime} = \sigma^{-n}(\mathfrak{a}_{n})\) 是递增的。由于环 A 是左诺特的，存在整数 \(m \geqslant 0\) ，使得 \(\mathfrak{a}_{n}^{\prime} = \mathfrak{a}_{n+1}^{\prime}\) 对 \(n \geqslant m\) 成立。由于

\(\sigma\) 是满射，我们有关系式

\[
\sigma (\mathfrak {a} _ {n}) = \mathfrak {a} _ {n + 1}\tag{7}
\]

对每个整数 \(n \geqslant m\) 成立。

设 \(\mathfrak{c}\) 是 B 的由 \(\mathfrak{b} \cap \mathrm{B}_m\) 生成的左理想。由于左 A-模 \(\mathrm{B}_m\) 是有限生成的且环 A 是左诺特的，左 A-模 \(\mathfrak{b} \cap \mathrm{B}_m\) 是有限生成的（VIII, p.~7, 命题 4 a)）。因此左理想 \(\mathfrak{c}\) 由 B 的一个有限子集生成。显然它含于 \(\mathfrak{b}\) 。我们来证明它等于 \(\mathfrak{b}\) ，办法是用归纳法证明对每个整数 \(n \geqslant 0\) 有

\[
\mathfrak {b} \cap \mathrm{B} _ {n} \subset \mathfrak {c}.\tag{8}
\]

关系式 (8) 对 \(n \leqslant m\) 按构造即成立。以下设 n 是一个整数 \(\geqslant m\) 且满足 \(b \cap B_{n} \subset c\) 。设 P 是 \(b \cap B_{n+1}\) 的一个元素。那么 \(\varphi_{n+1}(P)\) 属于 \(\mathfrak{a}_{n+1} = \sigma(\mathfrak{a}_{n})\) ，因而存在 \(b \cap B_{n}\) 的一个元素 Q，使得 \(\varphi_{n+1}(P) = \sigma(\varphi_{n}(Q))\) 。令 R = P - XQ。由关系式 (6)，我们有 \(\varphi_{n+1}(R) = 0\) ，即 \(R \in B_{n}\) 。由于 P 与 Q 都属于 B 的左理想 b，R 亦然；于是 R 与 Q 属于 \(b \cap B_{n}\) ，而由归纳假设它含于理想 c。因此 P 属于 c。~这就证明了 \(b \cap B_{n+1} \subset c\) 。

由此推出 \(\mathfrak{b}\) 等于 \(\mathfrak{c}\) ；因此它是 B 的一个有限生成理想。这就证明了环 B 是左诺特的。

如果环 A 的自同态 \(\sigma\) 不是自同构，那么环 \(A[X]_{\sigma,d}\) 不一定是左诺特的，即使 A 是诺特交换环时也是如此（VIII, p.~22, 习题 26）。

推论 1（希尔伯特）. --- 设 A 是诺特交换环。对每个整数 \(n \geqslant 0\) ，多项式代数 \(A[X_{1}, \ldots, X_{n}]\) 是诺特环。

这由定理 2 并顾及 III, §2, No.~9, p.~453 的命题 8 经归纳法得出。

推论 2. --- 设 A 是诺特交换环。由有限多个元素生成的交换 A-代数是诺特环。

这样的代数同构于形如 \(A[X_{1},\ldots,X_{n}]/\mathfrak{a}\) 的代数，其中 \(n \geqslant 0\) 且 a 是 \(A[X_{1},\ldots,X_{n}]\) 的一个理想。然后应用推论 1 以及 VIII, p.~7 的命题 5。

推论 3. --- 每个交换环都是一族右向有向的诺特子环的并。

事实上，设 A 是交换环。A 的由有限多个元素（作为 Z-代数）生成的子环由推论 2 知是诺特的。它们构成 A 的一族右向有向的子环，其并为 A。

